\documentclass[11pt,a4paper]{article}

% ── 中文支持 ──
\usepackage[UTF8]{ctex}
\setCJKmainfont{SimSun}[BoldFont=SimHei,ItalicFont=KaiTi]
\setCJKsansfont{SimHei}
\setCJKmonofont{FangSong}

% ── 数学与物理 ──
\usepackage{amsmath,amssymb,amsfonts}
\usepackage{mathtools}
\usepackage{physics}
\usepackage{siunitx}
\sisetup{detect-all, separate-uncertainty=true}

% ── 排版 ──
\usepackage[a4paper,margin=2.5cm]{geometry}
\usepackage{graphicx}
\usepackage{subcaption}
\usepackage{booktabs}
\usepackage{multirow}
\usepackage{enumitem}
\usepackage[titletoc]{appendix}

% ── 参考文献 ──
\usepackage[backend=biber,style=ieee,sorting=none]{biblatex}
\addbibresource{references.bib}

% ── 超链接 ──
\usepackage[colorlinks=true,
            linkcolor=blue,
            citecolor=blue,
            urlcolor=blue]{hyperref}

% ── 代码与提示框 ──
\usepackage{listings}
\usepackage{tcolorbox}

% ── 自定义命令 ──
\newcommand{\FT}{\mathcal{F}}
\newcommand{\IFT}{\mathcal{F}^{-1}}
\newcommand{\rep}{f_{\text{rep}}}
\newcommand{\ceo}{f_{\text{CEO}}}
\newcommand{\gvd}{\beta_2}
\newcommand{\spm}{\gamma}
\newcommand{\pdresp}{\mathcal{R}}
\newcommand{\Afield}{A}
\newcommand{\Atilde}{\widetilde{A}}

\title{\textbf{超快光脉冲的时域、光谱域与射频域信号分析}\
\large ——复光场振幅 $A(t)$ 是连接三个测量域的纽带}
\author{OUTERI}
\date{\today}

\begin{document}

\maketitle

\begin{abstract}
本文从超快光学中最基本却也最容易被混淆的一组关系出发——复光场振幅 $A(t)$ 与其傅里叶变换 $\Atilde(\omega)$
构成傅里叶对，而非光功率——系统阐述时域光功率、光学频谱（OSA）和射频谱（PD\,+\,
ESA）之间的数学联系与物理区别。核心结论为：
\[
\boxed{\FT\{|A(t)|^2\} \neq |\Atilde(\omega)|^2}
\]
前者（经平方律探测后的傅里叶变换）对应 RF 谱，反映光强的周期性和强度噪声；
后者（先傅里叶变换再取模平方）对应光学频谱，反映电场的频率分量。
以此为基础，本文依次讨论“时域脉冲”在仿真与实验中的具体含义、OSA
丢失谱相位所导致的脉宽推断局限、锁模光频梳的形成条件、PD
平方律探测将光频梳映射为 RF 谐波梳的严格推导、单脉冲到 ESA
射频谱的泊松求和表述、PD 示波器波形与实际飞秒脉冲的本质区别、ESA
能/不能回答的问题、以及主动锁模、孤子锁模、调Q 和调Q
锁模四种机制在三个测量域中的不同表现。最后结合色散管理片上激光器的典型参数，
说明 OSA、PD+ESA、自相关仪和光学外差测量各自回答的是不同物理问题。
\end{abstract}

\tableofcontents

% ======================================================================
\section{核心关系：$A(t) \longleftrightarrow \Atilde(\omega)$ 是傅里叶对}
% ======================================================================

超快光学测量中最基本、也最容易被混淆的一组关系是：

\begin{equation}\label{eq:core_fourier}
\boxed{A(t)\ \longleftrightarrow\ \Atilde(\omega)}
\end{equation}

互为傅里叶变换的是\textbf{复光场振幅} $A(t)$，不是光功率。
由此派生出三个实验可测量：

\begin{equation}\label{eq:three_domains}
\boxed{
\begin{aligned}
P(t) &= |A(t)|^2                     &&\text{时域光功率（PD 直流耦合示波器）} \[4pt]
S_{\rm opt}(\omega) &= |\Atilde(\omega)|^2  &&\text{光学频谱（OSA）} \[4pt]
S_{\rm RF}(f) &\propto \left|\FT\{P(t)\}\right|^2
                 = \left|\FT\{|A(t)|^2\}\right|^2  &&\text{射频谱（PD + ESA）}
\end{aligned}}
\end{equation}

这三个量的运算次序截然不同：

\begin{center}
\begin{tcolorbox}[width=0.92\textwidth, colback=blue!3, colframe=blue!40]
\textbf{三个域的运算次序：}
\begin{itemize}
    \item \textbf{OSA}：先 $\FT$，再 $|\cdot|^2$——测的是光场频谱模平方；
    \item \textbf{PD}：先 $|\cdot|^2$——测的是光场强度（功率），载波相位丢失；
    \item \textbf{ESA}：先 $|\cdot|^2$，再 $\FT$，再 $|\cdot|^2$——测的是光功率随时间变化所包含的射频频率分量。
\end{itemize}
\end{tcolorbox}
\end{center}

因此，光学频谱与射频谱不是同一个频谱，只是由同一光场 $A(t)$ 通过不同运算得到。
这一区分是理解锁模激光测量的关键。

% ======================================================================
\section{“时域脉冲”具体指什么}
% ======================================================================

实际电场可写为载波--包络形式：

\begin{equation}\label{eq:real_field}
E_{\rm real}(t)
= \operatorname{Re}\left\{A(t)\,e^{i\omega_0 t}\right\},
\end{equation}

其中：
\begin{itemize}
    \item $\omega_0 = 2\pi\nu_0$ 是光载波角频率；
    \item $A(t) = |A(t)|\,e^{i\phi(t)}$ 是\textbf{缓慢变化的复包络}（slowly-varying complex envelope）；
    \item 仿真和理论分析中所说的“时域脉冲”通常指 $A(t)$ 或 $|A(t)|^2$，不是飞快振荡的实际光载波。
\end{itemize}

在 $\SI{1550}{nm}$ 处：
\begin{equation}\label{eq:carrier_numbers}
\nu_0 = \frac{c}{\lambda} \approx \SI{193.4}{THz},
\qquad
T_0 = \frac{1}{\nu_0} \approx \SI{5.17}{fs}.
\end{equation}

一个 $\SI{300}{fs}$ 脉冲中大约包含几十个光学振荡周期。普通高速 PD 带宽只有数 GHz
至数十 GHz，完全跟不上 $\SI{193}{THz}$ 的光载波，因此只能测量光学周期平均后的功率：

\begin{equation}\label{eq:optical_power}
P(t) = \frac{n\varepsilon_0 c A_{\rm eff}}{2}\,|A(t)|^2.
\end{equation}

在光纤激光仿真中，常将 $A$ 归一化使 $|A(t)|^2 = P(t)$，单位直接取 W。

% ======================================================================
\section{光谱仪（OSA）测到的光学频谱}
% ======================================================================

\subsection{复包络的傅里叶变换}

复包络的傅里叶变换定义为：

\begin{equation}\label{eq:ft_envelope}
\Atilde(\Omega) = \int_{-\infty}^{\infty} A(t)\,e^{-i\Omega t}\,\dd{t},
\qquad \Omega = \omega - \omega_0.
\end{equation}

光谱仪测量的是\textbf{功率谱密度}：

\begin{equation}\label{eq:osa_measures}
S_{\rm opt}(\omega) \propto |\Atilde(\omega)|^2.
\end{equation}

\subsection{OSA 丢失了谱相位}

OSA 只保留傅里叶变换的幅度，完全丢失了谱相位：
\begin{equation}\label{eq:lost_phase}
\phi_{\rm spec}(\omega) = \arg \Atilde(\omega).
\end{equation}

因此，\textbf{仅凭 OSA 光谱不能唯一确定时域脉冲}。

\subsubsection{实例：两个光谱完全相同但时域不同的脉冲}

设有两个脉冲：
\begin{equation}\label{eq:two_pulses}
\begin{aligned}
\Atilde_1(\Omega) &= |\Atilde(\Omega)|, \[4pt]
\Atilde_2(\Omega) &= |\Atilde(\Omega)|\,e^{i\phi_2\Omega^2/2}.
\end{aligned}
\end{equation}

它们在 OSA 上\textbf{完全相同}，但第二个脉冲具有二阶谱相位（线性啁啾）。
对于高斯脉冲，实际脉宽与变换极限脉宽的关系为：

\begin{equation}\label{eq:chirped_width}
\tau = \tau_{\rm TL}\,
\sqrt{1 + \left(\frac{4\ln 2\,\phi_2}{\tau_{\rm TL}^2}\right)^2},
\end{equation}

其中 $\tau_{\rm TL}$ 是变换极限脉宽（$\phi_2 = 0$）。
即使光谱完全不变，$\phi_2$ 仍可使脉冲展宽数倍乃至数十倍。

\begin{tcolorbox}[width=0.92\textwidth, colback=red!3, colframe=red!50,
    title=重要结论]
\textbf{宽光谱不等于短脉冲。}
必须同时测量谱相位，或使用自相关（autocorrelator）、FROG、SPIDER、d-scan
等方法才能完整表征脉冲。
\end{tcolorbox}

% ======================================================================
\section{锁模脉冲为什么在光谱上表现为光频梳}
% ======================================================================

\subsection{纵模的相干叠加}

设腔内有一组纵模：
\begin{equation}\label{eq:longitudinal_modes}
E^{(+)}(t) = \sum_m a_m\,e^{i(2\pi\nu_m t + \phi_m)}.
\end{equation}

\textbf{锁模的本质要求}是不同纵模之间具有稳定的相位关系。
在理想锁模状态下：
\begin{equation}\label{eq:comb_freqs}
\nu_m = \ceo + m\rep,
\end{equation}

其中：
\begin{itemize}
    \item $\rep$：脉冲重复频率，通常等于腔自由光谱范围（FSR）；
    \item $\ceo$：载波--包络偏移频率（Carrier-Envelope Offset）；
    \item $m$：光梳线编号（$m \sim 10^5\text{--}10^6$）。
\end{itemize}

\subsection{相干叠加 vs.\ 非相干叠加}

若各纵模相位近似满足线性关系 $\phi_m = \phi_0 + m\phi_1$，
所有纵模在特定时刻相干叠加，形成周期性短脉冲。

设共有 $M$ 个等振幅纵模：
\begin{itemize}
    \item \textbf{相干叠加}（锁模）：峰值场 $\propto M$，峰值功率 $\propto M^2$；
    \item \textbf{非相干随机相位}（非锁模多纵模）：平均功率 $\propto M$。
\end{itemize}

所以“很多纵模同时振荡”不自动等于锁模；必须有稳定的纵模相位关系。
Haus 的经典综述系统给出了这一主方程框架 \cite{Haus2000}。

\subsection{载波--包络偏移频率 $\ceo$ 的物理起源}

由于群速度与相速度的差异，脉冲包络在每个往返中以群速度 $v_g$ 传播，
而载波以相速度 $v_p$ 传播。每经过一个腔往返，
载波相位相对于包络峰值滑移 $\Delta\phi_{\rm CEO}$：

\begin{equation}\label{eq:ceo_physics}
\ceo = \frac{\rep}{2\pi}\,\Delta\phi_{\rm CEO}
= \rep\left(\frac{L}{v_g} - \frac{L}{v_p}\right)\frac{\omega_0}{2\pi} \pmod{\rep}.
\end{equation}

$\ceo$ 的测量通常需要 $f\text{--}2f$ 等非线性自参考方法，
普通 PD + ESA 无法直接测量 \cite{Cundiff2003,Newbury2005}。

% ======================================================================
\section{PD 为什么把光频梳变成射频信号}
% ======================================================================

\subsection{平方律探测的严格推导}

理想线性光电探测器满足平方律。更实际的模型为：

\begin{equation}\label{eq:pd_model}
i(t) = \pdresp \int h_{\rm PD}(t-t')\,P(t')\,\dd{t'} + n(t),
\end{equation}

其中：
\begin{itemize}
    \item $\pdresp = \eta e/(h\nu_0)$：响应度（A/W），$\eta$ 为量子效率；
    \item $h_{\rm PD}(t)$：PD 的脉冲响应（由载流子渡越时间和 RC 时间常数决定）；
    \item $n(t)$：散粒噪声、热噪声等的总和。
\end{itemize}

将多纵模光场代入平方律（忽略 PD 带宽限制和噪声，先看理想情况）：

\begin{equation}\label{eq:pd_square_law}
\begin{aligned}
i(t) &\propto \left|\sum_m a_m\,e^{i2\pi\nu_m t}\right|^2 \[4pt]
&= \underbrace{\sum_m |a_m|^2}_{\text{直流项}}
   + \underbrace{\sum_{m\neq n} a_m a_n^*\,e^{i2\pi(\nu_m-\nu_n)t}}_{\text{拍频项}}.
\end{aligned}
\end{equation}

\subsection{拍频产生 RF 谐波}

因为 $\nu_m - \nu_n = (m-n)\rep$，所以 PD 输出包含：
\begin{equation}\label{eq:rf_harmonics}
\rep,\; 2\rep,\; 3\rep,\; \ldots
\end{equation}

对于第 $k$ 次射频谐波，其复振幅为：
\begin{equation}\label{eq:kth_harmonic}
I_k \propto \sum_m a_{m+k} a_m^*.
\end{equation}

\textbf{关键理解}：射频信号是所有相隔 $k$ 根光梳线之间拍频的\textbf{相干叠加}，
而不是某一根光梳线直接“变成”了射频。

\subsection{$\ceo$ 在直接光电探测中被消掉}

\begin{equation}\label{eq:ceo_cancellation}
(\ceo + (m+k)\rep) - (\ceo + m\rep) = k\rep.
\end{equation}

所以 $\ceo$ 在直接光电探测中消掉了。普通 PD + ESA 能测 $\rep$，
但一般不能直接测 $\ceo$；$\ceo$ 的测量通常需要 $f\text{--}2f$ 干涉等非线性自参考方法
\cite{Cundiff2003,Newbury2005}。

% ======================================================================
\section{从单脉冲到 ESA 射频谱的泊松求和推导}
% ======================================================================

设单个光功率脉冲为 $p(t)$，周期 $T_{\rm rep}$：

\begin{equation}\label{eq:pulse_train_power}
P(t) = \sum_{n=-\infty}^{\infty} p(t - nT_{\rm rep}).
\end{equation}

令 $\rep = 1/T_{\rm rep}$，$E_p = \int p(t)\,\dd{t}$（单脉冲能量）。
根据泊松求和公式：

\begin{equation}\label{eq:poisson_rf}
\widetilde{P}(f)
= \frac{1}{T_{\rm rep}} \sum_k \widetilde{p}(k\rep)\,\delta(f - k\rep).
\end{equation}

经过 PD 及微波链路后：

\begin{equation}\label{eq:esa_harmonics}
I(f) = \frac{\pdresp}{T_{\rm rep}} \sum_k H_{\rm PD}(k\rep)\,\widetilde{p}(k\rep)\,\delta(f - k\rep).
\end{equation}

因此 ESA 上第 $k$ 个谐波的功率满足：

\begin{equation}\label{eq:esa_power_k}
\boxed{P_{{\rm RF},k} \propto \pdresp^2\,|H_{\rm PD}(k\rep)|^2\,|\widetilde{p}(k\rep)|^2}.
\end{equation}

\textbf{这说明 ESA 谐波包络同时受两部分影响：}
\begin{enumerate}[label=(\arabic*)]
    \item 光功率脉冲形状 $\widetilde{p}(f)$（单脉冲强度傅里叶变换）；
    \item PD、电缆、放大器和 ESA 的总频率响应 $H_{\rm PD}(f)$。
\end{enumerate}

\subsection{飞秒脉冲的 ESA 谐波为何不能直接换算脉宽}

对于很短的飞秒脉冲，在电子学可测范围（通常至 $\sim\SI{50}{GHz}$）内通常满足：
\begin{equation}\label{eq:narrowband_rf}
k\rep\,\tau_p \ll 1,
\end{equation}

于是 $\widetilde{p}(k\rep) \approx E_p$（所有低次谐波幅度几乎相同）。
此时 ESA 上各谐波的衰减\textbf{主要反映 PD 和微波链路带宽}，而不是飞秒脉宽。

这就是为什么不能用“ESA 上看到了十次谐波”直接计算 $\SI{300}{fs}$ 脉宽。
直接探测超短脉冲列的严格噪声和射频推导可参考 Quinlan et al.\ \cite{Quinlan2013}。

% ======================================================================
\section{PD 示波器波形通常不是实际飞秒脉冲}
% ======================================================================

PD 的近似上升时间由带宽决定：
\begin{equation}\label{eq:pd_risetime}
t_r \approx \frac{0.35}{B_{\rm PD}}.
\end{equation}

\textbf{示例}：$B_{\rm PD} = \SI{12}{GHz}$ 的 PD，$t_r \approx \SI{29}{ps}$。
如果输入脉冲只有 $\SI{100}{fs}$，PD 输出不会呈现 $\SI{100}{fs}$ 宽度，
而是呈现约几十皮秒的电子脉冲响应：

\begin{equation}\label{eq:pd_output_fs}
i(t) \approx \pdresp\,E_p\,h_{\rm PD}(t).
\end{equation}

PD 输出脉冲的\textbf{面积}仍与光脉冲能量成正比（$\int i(t)\dd{t} \propto E_p$），
但输出波形\textbf{宽度}主要是探测器响应。

\begin{tcolorbox}[width=0.92\textwidth, colback=orange!3, colframe=orange!60,
    title=PD + 示波器能/不能做什么]
\textbf{能做：}
\begin{itemize}
    \item 确认脉冲周期 $\rep$ 和脉冲存在性；
    \item 观察调 Q 包络和脉冲幅度稳定性；
    \item 检测脉冲缺失（dropout）和多脉冲。
\end{itemize}
\textbf{不能做：}
\begin{itemize}
    \item 测量飞秒或皮秒脉宽；
    \item 飞秒脉宽需要自相关仪、FROG、SPIDER 等光学方法。
\end{itemize}
\end{tcolorbox}

% ======================================================================
\section{ESA 射频谱究竟能说明什么}
% ======================================================================

ESA 测量的是光电流（或微波电压）的功率谱密度：

\begin{equation}\label{eq:esa_definition}
S_i(f) = \lim_{T\to\infty} \frac{1}{T}\,|I_T(f)|^2.
\end{equation}

\subsection{ESA 能判断的信息}

稳定锁模时，ESA 出现窄的 $\rep$ 及其谐波峰。由此可以判断：
\begin{itemize}
    \item 重复频率 $\rep$（精确值）；
    \item 脉冲列是否周期稳定（谐波峰是否尖锐、有无分裂）；
    \item 是否存在低频调 Q 包络（kHz--MHz 区间的额外峰）；
    \item 是否有多脉冲或脉冲能量振荡（$\rep/2$ 等处出现亚谐波峰）；
    \item 重复频率相位噪声和定时抖动（从 $\rep$ 谐波的线宽和噪声基底推算）；
    \item 谐波锁模阶数（若 $\rep$ 峰微弱但 $q\rep$ 峰很强）。
\end{itemize}

\subsection{ESA 不能单独确定的信息}

\begin{itemize}
    \item 飞秒脉宽（需要光学自相关等方法）；
    \item 光学谱相位（ESA 无相位信息）；
    \item 是否达到变换极限；
    \item 每根光梳线的绝对光学线宽（需要光学外差测量）；
    \item $\ceo$（需要 $f\text{--}2f$ 干涉等非线性方法）。
\end{itemize}

\subsection{从 RF 相位噪声推算定时抖动}

如果重复频率信号写为：
\begin{equation}\label{eq:rf_with_phase_noise}
V(t) = V_0\cos\!\big[2\pi\rep t + \phi(t)\big],
\end{equation}

则相位波动 $\phi(t)$ 对应的时间误差为：
\begin{equation}\label{eq:timing_error}
\delta t(t) = \frac{\phi(t)}{2\pi\rep}.
\end{equation}

由单边带相位噪声 $L(f)$（单位：dBc/Hz）可计算积分定时抖动：

\begin{equation}\label{eq:integrated_jitter}
\boxed{\sigma_t
= \frac{1}{2\pi\rep}\,
\sqrt{2\int_{f_1}^{f_2} 10^{L(f)/10}\,\dd{f}} }.
\end{equation}

经典射频噪声分析参见 von der Linde \cite{vonderLinde1986} 和 Paschotta \cite{Paschotta2004}。

% ======================================================================
\section{三种（加一种）脉冲产生机制的对比分析}
% ======================================================================

主动锁模、孤子锁模和调Q 脉冲\textbf{不是同一分类层级}的概念：

\begin{table}[htbp]
\centering
\caption{四种脉冲产生机制的分类依据与特征对比}
\label{tab:classification}
\renewcommand{\arraystretch}{1.5}
\begin{tabular}{@{}p{2cm} p{2.2cm} p{2.8cm} p{2.2cm} p{2.4cm}@{}}
\toprule
\textbf{类型} & \textbf{分类依据} & \textbf{核心机制} & \textbf{典型时间尺度} & \textbf{ESA 主要特征} \
\midrule
主动锁模 & 锁模方式 & 外部调制器周期性调制损耗或相位 & ps--ns，MHz--GHz & 强窄线位于调制频率及谐波 \
\hline
孤子锁模 & 脉冲整形机制 & 反常色散与 Kerr 非线性平衡 & fs--ps，MHz--GHz & 清洁的 $\rep$ 谐波列 \
\hline
调Q & 储能释放机制 & 先低 Q 储能，再高 Q 快速释放 & ns--$\mu$s，Hz--MHz & 主要在 $f_Q$ 及其谐波 \
\hline
调Q 锁模 & 两种动力学共存 & 调Q 包络中包含锁模子脉冲 & 慢包络 + 快脉冲列 & $f_Q$、$\rep$ 同时出现，$\rep$ 周围有 $f_Q$ 边带 \
\bottomrule
\end{tabular}
\end{table}

% ======================================================================
\section{主动锁模}
% ======================================================================

\subsection{物理机制}

主动锁模通过电光或声光调制器在腔内产生周期性损耗窗口：

\begin{equation}\label{eq:am_modulation}
\ell(t) = \ell_0 + \Delta\ell\big[1 - \cos(\Omega_m t)\big].
\end{equation}

当调制频率 $\Omega_m \approx q\,\Omega_{\rm FSR}$ 时（$q=1$ 基频锁模，
$q>1$ 谐波锁模），脉冲每次经过调制器时都落在低损耗窗口中。

\subsection{基本特征}

\begin{itemize}
    \item 重复频率主要由外部射频源和腔长共同决定；
    \item RF 线通常与驱动源高度相关，相位噪声受限于微波源；
    \item 定时抖动可被低噪声微波源压低；
    \item 脉冲通常近似高斯形，但并非必然是孤子；
    \item 调制频率失谐会引起脉冲漂移、展宽或超模噪声。
\end{itemize}

理论基础参见 Kuizenga \& Siegman \cite{Kuizenga1970}。

\subsection{数值例子：$\SI{10}{GHz}$ 主动锁模}

假设：$\rep = \SI{10}{GHz}$，$P_{\rm avg} = \SI{20}{mW}$，$\tau_{\rm FWHM} = \SI{2}{ps}$。

单脉冲能量：
\begin{equation}\label{eq:am_example_energy}
E_p = \frac{P_{\rm avg}}{\rep} = \frac{\SI{20}{mW}}{\SI{10}{GHz}} = \SI{2}{pJ}.
\end{equation}

高斯脉冲 $E_p = 1.0645\,P_0\,\tau_{\rm FWHM}$，峰值功率：
\begin{equation}\label{eq:am_example_peak}
P_0 \approx \frac{\SI{2}{pJ}}{1.0645 \times \SI{2}{ps}} \approx \SI{0.94}{W}.
\end{equation}

若为变换极限高斯脉冲（$\Delta\nu \cdot \tau_{\rm FWHM} = 0.441$）：
\begin{equation}\label{eq:am_example_spectrum}
\Delta\nu \approx \frac{0.441}{\SI{2}{ps}} = \SI{220.5}{GHz},
\qquad
\Delta\lambda \approx \SI{1.77}{nm}\;(\text{@ }\SI{1550}{nm}).
\end{equation}

$\SI{12}{GHz}$ PD 通常只能响应直流和 $\SI{10}{GHz}$ 基频，
示波器上可能接近正弦波，而不是清晰的 $\SI{2}{ps}$ 脉冲。

% ======================================================================
\section{孤子锁模}
% ======================================================================

\subsection{NLSE 与基态孤子解}

在忽略增益、损耗和滤波的理想反常色散介质中，
脉冲演化由非线性薛定谔方程（NLSE）描述：

\begin{equation}\label{eq:nlse_soliton}
\frac{\partial A}{\partial z}
= -i\frac{\gvd}{2}\frac{\partial^2 A}{\partial t^2}
+ i\spm|A|^2 A,
\qquad \gvd < 0.
\end{equation}

基态孤子解为：
\begin{equation}\label{eq:fundamental_soliton}
A(z,t) = \sqrt{P_0}\,\sech\!\left(\frac{t}{T_0}\right)e^{iz/(2L_D)},
\end{equation}

孤子阶数 $N^2 = \gamma P_0 T_0^2/|\gvd| = 1$（基态孤子条件）。
由此：
\begin{equation}\label{eq:soliton_relations}
P_0 = \frac{|\gvd|}{\spm T_0^2},
\qquad
E_p = 2P_0 T_0,
\qquad
\tau_{\rm FWHM} = 2\ln(1+\sqrt{2})\,T_0 \approx 1.763\,T_0.
\end{equation}

变换极限 $\sech^2$ 脉冲的时间--带宽积：
\begin{equation}\label{eq:soliton_tbp}
\Delta\nu \cdot \tau_{\rm FWHM} = 0.315.
\end{equation}

\subsection{Haus 主方程与 Kelly 边带}

实际锁模激光器是耗散系统，需用 Haus 主方程描述增益、损耗、饱和吸收和滤波的平衡
\cite{Haus2000,Ippen1994}。孤子仅是色散与非线性平衡形成的脉冲整形状态。
色散管理孤子的更全面综述见 Turitsyn et al.\ \cite{Turitsyn2012}。

周期性腔元件对孤子产生扰动时，光谱上可能出现 \textbf{Kelly 边带} \cite{Kelly1992}。
但“没有明显 Kelly 边带”不等于不是孤子——色散管理腔中的相位匹配条件可能受到抑制。

\subsection{数值例子：$\SI{100}{fs}$ 孤子脉冲}

假设：$\rep = \SI{100}{MHz}$，$P_{\rm avg} = \SI{10}{mW}$，$\tau_{\rm FWHM} = \SI{100}{fs}$。

\begin{equation}\label{eq:soliton_example}
\begin{aligned}
E_p &= \frac{\SI{10}{mW}}{\SI{100}{MHz}} = \SI{100}{pJ}, \[4pt]
T_0 &= \frac{\SI{100}{fs}}{1.763} \approx \SI{56.7}{fs}, \[4pt]
P_0 &= \frac{E_p}{2T_0} \approx \SI{881}{W}.
\end{aligned}
\end{equation}

变换极限光谱宽度：
\begin{equation}\label{eq:soliton_example_spectrum}
\Delta\nu = \frac{0.315}{\SI{100}{fs}} = \SI{3.15}{THz},
\qquad
\Delta\lambda \approx \SI{25.2}{nm}\;(\text{@ }\SI{1550}{nm}).
\end{equation}

对于 $\SI{12}{GHz}$ PD：
\begin{itemize}
    \item 示波器只能看到约 $\SI{29}{ps}$ 的电子响应，完全无法分辨 $\SI{100}{fs}$ 脉宽；
    \item ESA 可看见 $\SI{100}{MHz}$、$\SI{200}{MHz}$、$\SI{300}{MHz}\ldots$ 谐波；
    \item 这些谐波\textbf{不能}直接重构 $\SI{100}{fs}$ 光脉冲——它们受限于 PD 带宽，不是脉宽的量度。
\end{itemize}

% ======================================================================
\section{调Q 脉冲}
% ======================================================================

\subsection{物理机制：速率方程}

调Q 的核心不是纵模相位锁定，而是通过调制腔损耗来控制储能与释放。
反转粒子数 $N$ 和腔内光子数 $S$ 的耦合速率方程为：

\begin{equation}\label{eq:qswitch_rate_equations}
\begin{aligned}
\frac{\dd{N}}{\dd{t}} &= R_p - \frac{N}{\tau_f} - \sigma_e v_g N S, \[4pt]
\frac{\dd{S}}{\dd{t}} &= \left[\Gamma\sigma_e v_g N - \frac{1}{\tau_p(t)}\right]S
                        + \beta_{\rm sp}\frac{N}{\tau_f}.
\end{aligned}
\end{equation}

低 Q 状态下，$\tau_p$ 小、腔损耗大，光子不能积累但反转粒子数持续升高。
突然提高 Q 后，阈值下降，光子数快速增长并耗尽反转，形成高能量巨脉冲。

经典参考：McClung \& Hellwarth \cite{McClung1962}，Degnan \cite{Degnan1995}。

\subsection{调Q 脉冲的 RF 谱特征——与锁模的本质区别}

调Q 激光器的各纵模之间\textbf{无固定相位关系}。虽然每个纵模对 $(m, m+k)$
仍会产生频率为 $k\cdot\Delta\nu_{\rm FSR}$ 的拍频信号，
但由于各纵模的相位 $\phi_m$ 随机且统计独立：

\begin{equation}\label{eq:qswitch_rf_incoherent}
\begin{aligned}
\FT[i(t)](f) &\propto \sum_{m,n} |a_m||a_n|\,\delta\big(f - |\nu_m-\nu_n|\big)\,e^{i(\phi_m-\phi_n)}, \[4pt]
\mathbb{E}\big[|i(f)|^2\big] &\propto \sum_{m,n} |a_m|^2|a_n|^2
                              \approx N^2 \cdot \overline{|a|^4}.
\end{aligned}
\end{equation}

结果是 RF 谱为\textbf{平滑的宽带连续谱}，不形成离散频率梳。
这与锁模激光器的锐利 RF 谐波梳形成鲜明对比。

\subsection{数值例子：$\SI{20}{ns}$ 调Q 激光}

假设：$f_Q = \SI{20}{kHz}$，$P_{\rm avg} = \SI{1}{W}$，$\tau_Q = \SI{20}{ns}$。

\begin{equation}\label{eq:qswitch_example}
\begin{aligned}
E_Q &= \frac{\SI{1}{W}}{\SI{20}{kHz}} = \SI{50}{\mu J}, \[4pt]
P_0 &\approx \frac{E_Q}{1.0645\,\tau_Q} \approx \SI{2.35}{kW}.
\end{aligned}
\end{equation}

ESA 主要看到 $\SI{20}{kHz}$、$\SI{40}{kHz}$、$\SI{60}{kHz}\ldots$，
\textbf{而不是}由腔长决定的数百 MHz 或 GHz 锁模重复频率。

$\SI{20}{ns}$ 高斯脉冲的变换极限带宽仅约 $\Delta\nu_{\rm TL} \approx \SI{22}{MHz}$。
但实际调Q 激光可能因多纵模同时振荡、动态频率啁啾和空间烧孔效应
表现出远超 $\SI{22}{MHz}$ 的 OSA 光谱宽度。
\textbf{这个光谱宽度不能像理想锁模脉冲那样直接换算脉宽}。

% ======================================================================
\section{调Q 锁模：最容易误判的情况}
% ======================================================================

\subsection{数学模型}

调Q 锁模（Q-switched mode-locking, QML）的光功率可近似写为：

\begin{equation}\label{eq:qml_power}
P(t) = Q(t)\,\sum_n p(t - nT_{\rm rep}),
\end{equation}

其中 $Q(t)$ 是缓慢的调Q 包络（ns--$\mu$s 尺度），
$p(t)$ 是单个锁模子脉冲（fs--ps 尺度）。

\subsection{频域特征}

频域中乘积对应卷积，因此在 ESA 上：
\begin{itemize}
    \item 低频区出现 $f_Q$ 及其谐波；
    \item 高频区出现 $\rep$；
    \item 每个 $k\rep$ 周围出现\textbf{间隔为 $f_Q$ 的边带}；
    \item $\rep$ 射频峰可能仍然很高（$> \SI{60}{dB}$ 信噪比），
    但脉冲能量在慢时间尺度上剧烈振荡。
\end{itemize}

\subsection{如何避免误判}

仅报告“在 $\rep$ 处有一个 $\SI{60}{dB}$ 信噪比射频峰”仍不足以排除调Q 锁模。
还需完成以下检查：

\begin{enumerate}[label=(\arabic*)]
    \item 将 ESA 起始频率降至接近 DC，检查 kHz--MHz 低频区是否有额外峰；
    \item 用长时间窗示波器（$\sim\mu\text{s--ms}$）查看是否存在脉冲包络调制；
    \item 检查 $\rep$ 峰周围是否存在规则的 $f_Q$ 边带；
    \item 测量逐脉冲能量或 RF 相位噪声——QML 的相位噪声通常显著高于 CW 锁模。
\end{enumerate}

% ======================================================================
\section{实例分析：色散管理片上锁模激光器}
% ======================================================================

以典型色散管理片上锁模激光器的参数为例，说明各测量量之间的关系。

\subsection{已知参数}

\begin{itemize}
    \item $\rep = \SI{1.2}{GHz}$；
    \item ESA 上分辨出超过十次 $\rep$ 谐波；
    \item 光谱 3-dB 带宽 $\Delta\nu = \SI{1.14}{THz}$（$\Delta\lambda \approx \SI{9.5}{nm}$ @ $\SI{1550}{nm}$）；
    \item 约包含 950 根光梳线；
    \item 时域模拟呈 $\sech^2$ 轮廓；
    \item 光栅处啁啾换符号（色散管理）；
    \item 未观察到调Q 不稳定性。
\end{itemize}

\subsection{梳齿数量的验证}

\begin{equation}\label{eq:comb_line_count}
\frac{\Delta\nu}{\rep} = \frac{\SI{1.14}{THz}}{\SI{1.2}{GHz}} \approx 950.
\end{equation}

PD 中相邻梳齿两两拍频，全部贡献到 $\SI{1.2}{GHz}$ 基频；
相隔两根的梳齿贡献到 $\SI{2.4}{GHz}$，以此类推。
因此 ESA 上出现十多个谐波是自然结果。

\subsection{三组测量回答不同的问题}

\begin{table}[htbp]
\centering
\caption{不同测量手段回答的物理问题}
\label{tab:measurements}
\renewcommand{\arraystretch}{1.5}
\begin{tabular}{@{}p{3.2cm} p{5.8cm} p{5.5cm}@{}}
\toprule
\textbf{测量手段} & \textbf{回答的问题} & \textbf{注意事项} \
\midrule
OSA & 光谱有多宽？有多少根梳齿？ & 不能单独确定脉宽 \
\hline
PD + ESA & 脉冲列是否以 $\SI{1.2}{GHz}$ 周期稳定重复？ & 不能测量飞秒脉宽 \
\hline
自相关仪/FROG & 脉冲究竟有多短？是否达到变换极限？ & 需先经色散补偿 \
\hline
$\rep$ 相位噪声 & 脉冲到达时间有多稳定？（定时抖动） & 反映脉冲列整体稳定性 \
\hline
光学外差 & 单根光梳线的光学线宽和频率漂移 & 同时包含 $\rep$ 和 $\ceo$ 的噪声贡献 \
\bottomrule
\end{tabular}
\end{table}

\subsection{两种“线宽”不可混淆}

\begin{itemize}
    \item \textbf{重复频率 RF 线宽}：描述脉冲列定时稳定性，由 PD + ESA
    在 $\rep$ 处直接测量；
    \item \textbf{单根光梳线的光学线宽}：描述单频光的相干性，
    需用窄线宽 CW 参考激光进行光学外差测量，结果同时包含 $\rep$、
    $\ceo$ 及其他腔噪声的贡献。
\end{itemize}

% ======================================================================
\section{总结}
% ======================================================================

本文的核心信息可凝练为以下一组方程：

\begin{equation}\label{eq:final_summary}
\boxed{
\begin{aligned}
A(t) &\longleftrightarrow \Atilde(\omega)
      &&\text{复光场傅里叶对——一切分析的出发点} \[4pt]
|A(t)|^2 &= P(t)
         &&\text{时域光功率（PD 直流输出）} \[4pt]
|\Atilde(\omega)|^2 &= S_{\rm opt}(\omega)
                      &&\text{OSA 光谱——先 $\FT$，再 $|\cdot|^2$} \[4pt]
\big|\FT\{P(t)\}\big|^2 &= S_{\rm RF}(f)
                          &&\text{PD + ESA 射频谱——先 $|\cdot|^2$，再 $\FT$，再 $|\cdot|^2$}
\end{aligned}}
\end{equation}

其中最容易被忽略的核心不等式为：

\begin{equation}\label{eq:key_inequality}
\boxed{\FT\{|A(t)|^2\} \neq |\Atilde(\omega)|^2}
\end{equation}

前者（RF 谱）是光谱场的频差自相关，反映重复频率和强度噪声；
后者才是光学频谱。理解这一区分，是正确判读锁模激光测量数据的前提。

% ======================================================================
% 参考文献
% ======================================================================
\printbibliography[heading=bibintoc, title={参考文献}]

% ======================================================================
% 附录
% ======================================================================
\appendix

\section{泊松求和公式的推导}

周期脉冲串 $P(t) = \sum_n p(t - nT_{\rm rep})$ 的傅里叶变换：

\begin{equation}
\begin{aligned}
\widetilde{P}(f) &= \int_{-\infty}^{\infty} \sum_n p(t - nT_{\rm rep})\,e^{-i2\pi f t}\,\dd{t} \
&= \sum_n e^{-i2\pi f n T_{\rm rep}} \int_{-\infty}^{\infty} p(t')\,e^{-i2\pi f t'}\,\dd{t'} \
&= \widetilde{p}(f) \sum_n e^{-i2\pi f n T_{\rm rep}}.
\end{aligned}
\end{equation}

由泊松求和公式：$\displaystyle\sum_n e^{-i2\pi f n T_{\rm rep}} = \frac{1}{T_{\rm rep}}\sum_k \delta(f - k/T_{\rm rep})$。
代入即得式 (\ref{eq:poisson_rf})。

\section{定时抖动公式的推导}

设重复频率信号 $V(t) = V_0\cos[2\pi\rep t + \phi(t)]$。
相位噪声单边带功率谱 $L(f)$（单位：dBc/Hz）定义为
$L(f) = 10\log_{10}[S_\phi(f)/2]$，其中 $S_\phi(f)$ 是 $\phi(t)$ 的双边带功率谱密度。

均方相位抖动：
$\langle\phi^2\rangle = 2\int_{f_1}^{f_2} S_\phi(f)\,\dd{f}
                     = 2\int_{f_1}^{f_2} 2\cdot 10^{L(f)/10}\,\dd{f}$。

均方时间抖动：
$\sigma_t^2 = \langle\phi^2\rangle/(2\pi\rep)^2$，即式 (\ref{eq:integrated_jitter})。

\end{document}
